\section{8.1 Single-session preference (SSP): a v7 diagnostic
taxonomy}\label{single-session-preference-ssp-a-v7-diagnostic-taxonomy}

In the contaminated v7 run, SSP is the controller's weakest competency:
\textbf{11/30 correct (0.3667), CI {[}0.20, 0.5333{]}}. The failure-only
replay repaired 0/19 of those v7 failures, but it used different prompt
contexts and therefore does not isolate an answer-model effect. The
complementary top-up observes \textbf{8/29 correct} with one unresolved
item (maximum \textbf{9/30}), proving that the v7 count cannot be
carried forward as a cross-version result. The taxonomy below is
consequently a causal analysis of the \textbf{19 v7 failures}, not a
final-1.1 evaluation
(\texttt{docs/research/SINGLE\_SESSION\_PREFERENCE\_FAILURE\_ANALYSIS.md}).

\textbf{Primary failure categories (n = 19):}

{\def\LTcaptype{none} % do not increment counter
\begin{longtable}[]{@{}lll@{}}
\toprule\noalign{}
Primary category & Count & Share \\
\midrule\noalign{}
\endhead
\bottomrule\noalign{}
\endlastfoot
retrieval-omission & 6 & 31.6\% \\
cross-session-contamination & 4 & 21.1\% \\
target-attribute-extraction-miss & 3 & 15.8\% \\
distractor-entity-binding & 2 & 10.5\% \\
polarity-error & 2 & 10.5\% \\
unsupported-answer-specificity & 1 & 5.3\% \\
immediate-request-vs-enduring-preference & 1 & 5.3\% \\
\end{longtable}
}

\textbf{Combined mechanism mentions (primary + secondary, of 19):}
retrieval-omission 11; cross-session-contamination 8;
distractor-entity-binding 7; target-attribute-extraction-miss 5;
unsupported-answer-specificity 3;
immediate-request-vs-enduring-preference 2; polarity-error 2.

\textbf{Pipeline-layer localization:} retrieval-ranking 12/19 (63.2\%);
answer-realization 10/19 (52.6\%); context-construction (top-k trim)
5/19 (26.3\%); attribution/target-binding 5/19 (26.3\%); ingestion,
extraction, persistence, controller-policy, and judge-scorer 0/19 each.
Notably, \textbf{no SSP failure was traced to the write-governance
policy itself}: the decisive facts exist in the corpus, and for the
answer-side cases they were even present in the prompt.

Two families dominate. \textbf{Grounding-supply failures}
(retrieval-omission + cross-session-contamination: primary in 10/19,
present in 15/19): the single supporting session's decisive preference
turn is out-ranked by topically similar distractor sessions and trimmed
away, and the model then answers from the wrong session or fabricates.
The 0/19 replay repair is consistent with that localization but is not
causal proof because the replay changed prompt contexts as well as the
answer path. \textbf{Answer-side binding/realization failures} (primary
in 9/19): the decisive fact was in context, but the answer bound the
preference to the wrong target, violated an explicitly stated negative
constraint (recommending a dis-preferred podcast genre), or
over-specified beyond evidence. One item (\texttt{1da05512}) is a
borderline judge-severity case rather than a clean pipeline defect, and
several judge rationales say ``fabricates'' where the mechanism is
distractor binding --- the offending tokens are verbatim in retrieved
memory.

\textbf{Highest-leverage proposed fix (hypothesized, not run):} make
preference grounding an explicit, provenance-constrained representation
--- a typed record (holder / target / attribute / polarity / scope /
source-session) whose retrieval is guaranteed for the question's subject
--- paired with an answer contract that asserts only
retrieved-and-grounded preferences and honors stated negatives. The
design targets the \textasciitilde15/19 grounding-supply mentions, while
the \textasciitilde9/19 answer-side failures show why a retrieval-only
repair is unlikely to suffice. Necessity and effect size remain
untested. The design and its credit-attribution ablation matrix are
specified in
\texttt{docs/research/SINGLE\_SESSION\_PREFERENCE\_METHOD.md};
PrefEval's finding that preference following collapses within ten turns
\citep{zhao2025prefeval} independently motivates the same construct.

\section{8.2 MGP-6 error profiles: three systems, three
signatures}\label{mgp-6-error-profiles-three-systems-three-signatures}

\textbf{Reference policy --- conservative, one-directional on the full
pilot.} Its 14 full-corpus action errors are SUPERSEDE→ABSTAIN. On test
it has 0/33 action false commits, but the one-sided 95\% upper bound is
0.0868 and it misses 10/45 updates. This is a measured trade-off inside
the pilot, not a population-safety result.

\textbf{Structured LLM policy --- different error allocation.} Test
action accuracy is 0.9000 versus 0.8889 for the reference, with a paired
family interval containing zero. It has more action false commits (2/33)
but fewer missed updates (7/45). This supports cost-sensitive reporting,
not an absolute policy ordering.

\textbf{Zero-shot decision baseline --- diagnostic gap.} Its lower score
and high full-corpus invalid-commit rate show a failure mode worth
studying, but label leakage prevents attributing the difference cleanly
to the governance scaffold.

\textbf{Ablation signatures.} Abstention removal → invalid-commit 1.0;
typing removal → KEEP F1 0.0; lineage removal → audit-score decay
invisible to action accuracy (0.8556→0.3689). These are distinct
within-generator failure signatures, consistent with the motivation
behind MemGuard's typing argument \citep{ha2026memguard} and TOKI's
audit rows \citep{wang2026toki}. They motivate independent component
tests; they do not by themselves establish necessity, causal effect
size, or external validity.

\section{8.3 LongMemEval run-integrity failure classes
(F1--F12)}\label{longmemeval-run-integrity-failure-classes-f1f12}

The failure ledger
(\texttt{docs/audit/CONTROLLERRUN\_FAILURE\_LEDGER.md}) distinguishes
per-item failures (§8.1) from run/pipeline failures, of which the most
consequential are: \textbf{F1/F2} --- HTTP-402 provider no-responses
cached as valid empty extractions (316 ``402'' traces; 388 empty
extractions), fixed in v8 by fail-closed extraction; \textbf{F9} ---
relation-classifier no-responses converted to a cached semantic Keep,
same fix; \textbf{F3/F4} --- total-loss replay failures (thinking-mode
token exhaustion → 0 predictions; embedder segfaults), fixed by
\texttt{enable\_thinking:false} and a shared embedding server;
\textbf{F6} --- judge false negatives on number-words and dates (``4''
vs ``four'', ``Valentine's Day'' vs ``February 14th''), fixed by guarded
normalization with deferral to the semantic judge; \textbf{F8} --- the
v7 escalation no-op; \textbf{F5/F11} --- provenance defects (misnamed
aborted run; a manifest pointing to a non-existent replay directory);
\textbf{F10} --- in-place verdict overrides preserved with originals;
and the top-up evidence showing that prior successes cannot be retained
by assumption. The general lesson: infrastructure failures that default
to ``empty but valid'' silently corrupt memory state, while selective
replay silently biases reported accuracy; fallbacks must fail closed and
repairs must be evaluated on homogeneous coverage.

\section{8.4 The retained annotation disagreement as construct
evidence}\label{the-retained-annotation-disagreement-as-construct-evidence}

\texttt{MGP6V2-T3-TEST-016} (implicit conflict): gold SUPERSEDE,
independent annotator ABSTAIN. The disagreement sits on the same
SUPERSEDE/ABSTAIN boundary as the reference policy's conservative
errors, making it a useful candidate for independent adjudication. With
one human label and an author-aligned manual, it is not enough to
identify a generally hard construct region. It is retained for blinded
adjudication; sensitivity is bounded (one item = 1.1\% of the test
split) and both labels are defensible under the current manual.

\section{8.5 Decision--state--answer
divergence}\label{decisionstateanswer-divergence}

The pilot illustrates why the three layers should be reported
separately. Within MGP-6, similar action accuracy can coexist with very
different premise-resistance behavior, and the no-lineage ablation
leaves action labels unchanged while changing the instrument's state
metric. Separately, the contaminated LongMemEval slice happens to rank
B4 naive overwrite above the controller on answers (63.3\% vs 61.1\%).
Because the first observations are author-aligned and the second is
contaminated, these are diagnostic examples rather than clean
cross-layer effect estimates. They nevertheless show how a single
aggregate can conceal different failure locations and motivate the
layered evaluation design (RQ3).

\section{8.6 Asymmetric harm and the price of
conservatism}\label{asymmetric-harm-and-the-price-of-conservatism}

The reference policy's full-pilot action errors are conservative, but
missed supersession can itself be harmful. Whether the trade is right is
policy-dependent: under the frozen commit-averse matrix the reference
has lower test loss (0.111 vs 0.189), while under the update-responsive
matrix the structured LLM has lower loss (0.433 vs 0.556). The current
weights and any prototype score thresholds have not been validated on
independent data, so this is sensitivity analysis rather than
optimization evidence.

\section{8.7 What the near-neighbor cluster
implies}\label{what-the-near-neighbor-cluster-implies}

Five verified 2026-adjacent works (STALE, Supersede, TOKI, MemGuard,
HaluMem) cover read-side staleness probing, the supersession gap, formal
write-time contradiction operators, write-time typing, and
operation-level evaluation. This thesis is positioned as an
\textbf{evaluation decomposition and validity audit}: it makes no claim
that the individual mechanisms or their conjunction are unprecedented.
The convergence supports the importance of the problem, while the
deterministic-policy finding of \citep{reddy2026freshness} motivates
independent confirmation rather than validating this pilot by analogy.
